%==============================================================
% PROJECT 14 — U-NET SEMANTIC LANE SEGMENTATION
% BRAND-NEW 6-PAGE AUTONOMY REDESIGN
%==============================================================
%
% DESIGN GOALS
% -------------------------------------------------------------
% • Explain semantic segmentation to a non-ML engineer before showing metrics.
% • Never reuse an image.
% • Keep annotation quality, class histograms, IoU, loss, and predictions
%   next to the evidence they explain.
% • Avoid dense LaTeX tables; use simple visual scorecards instead.
% • Match the modern header/footer language used by Projects 11–13.
%
% REQUIRED IMAGE FILES — EACH USED EXACTLY ONCE
% -------------------------------------------------------------
% behind_shot.jpg
% intanceType.png
% histogram1.png
% annotationMask1.png
% lossANDiou1.png
% histogram2.png
% lossANDiou2.png
% lossANDiou3.png
% lossANDiou4.png
% annotationMask2.png
% annotationMask3.png
% annotationMask4.png
% physical_framework.jpg
%
% Assumes the portfolio already defines:
% DeepGreen, PolyGold, SoftGreen, SoftGold, Muted, Ink, RuleGray
% and loads TikZ, graphicx, hyperref, fontawesome5.
%==============================================================

\newcommand{\UNPageBase}{%
    \fill[white] (current page.south west) rectangle (current page.north east);
    \fill[DeepGreen] (current page.north west) rectangle ($(current page.north east)+(0,-0.095in)$);
    \fill[PolyGold] ($(current page.north west)+(0,-0.095in)$) rectangle ($(current page.north east)+(0,-0.14in)$);
    \fill[DeepGreen] (current page.south west) rectangle ($(current page.south east)+(0,0.36in)$);
    \fill[PolyGold] ($(current page.south west)+(0,0.36in)$) rectangle ($(current page.south east)+(0,0.405in)$);
}

\newcommand{\UNHeader}[3]{%
    \node[anchor=north west,text=PolyGold,font=\sffamily\bfseries\fontsize{8.1}{8.7}\selectfont]
    at ($(current page.north west)+(0.42in,-0.39in)$)
    {03 AUTONOMY \hspace{5pt}/\hspace{5pt} PROJECT 14
    \ifx&#1&\else \hspace{5pt}/\hspace{5pt} #1 \fi};

    \node[anchor=north west,text=DeepGreen,font=\bfseries\fontsize{21.0}{22.0}\selectfont]
    at ($(current page.north west)+(0.40in,-0.67in)$) {#2};

    \node[anchor=north west,text=Muted,font=\sffamily\fontsize{8.15}{8.95}\selectfont]
    at ($(current page.north west)+(0.42in,-1.12in)$) {#3};

    \fill[PolyGold]
    ($(current page.north west)+(0.42in,-1.36in)$)
    rectangle ($(current page.north east)+(-0.42in,-1.315in)$);
}

\newcommand{\UNFooter}[3]{%
    \node[anchor=south west,text=white,font=\sffamily\bfseries\fontsize{7.0}{7.6}\selectfont]
    at ($(current page.south west)+(0.42in,0.17in)$) {#1};
    \node[anchor=south,text=white,font=\sffamily\fontsize{7.0}{7.6}\selectfont]
    at ($(current page.south)+(0,0.17in)$) {14 U-NET LANE SEGMENTATION \hspace{6pt} #2};
    \node[anchor=south east,text=white,font=\sffamily\bfseries\fontsize{7.0}{7.6}\selectfont]
    at ($(current page.south east)+(-0.42in,0.17in)$) {#3};
}

\newcommand{\UNMetric}[6]{%
    \fill[SoftGreen,rounded corners=3pt]
    ($(current page.north west)+(#1,#2)$) rectangle ++(#3,-0.78in);
    \node[anchor=north west,text=DeepGreen,font=\sffamily\bfseries\fontsize{6.8}{7.4}\selectfont]
    at ($(current page.north west)+(#1+0.16in,#2-0.12in)$) {#4};
    \node[anchor=north west,text=Ink,font=\bfseries\fontsize{13.0}{13.5}\selectfont]
    at ($(current page.north west)+(#1+0.16in,#2-0.36in)$) {#5};
    \node[anchor=north west,text=Muted,font=\sffamily\fontsize{5.9}{6.5}\selectfont]
    at ($(current page.north west)+(#1+0.16in,#2-0.61in)$) {#6};
}

%==============================================================
% PAGE 1 — WHAT SEMANTIC SEGMENTATION IS
%==============================================================
\clearpage
\thispagestyle{empty}
\hypertarget{unet}{}
\null
\begin{tikzpicture}[remember picture,overlay]
\UNPageBase
\UNHeader
{}
{U-NET SEMANTIC LANE SEGMENTATION}
{Cal Poly Autonomous Vehicles Laboratory \hspace{7pt}\textbullet\hspace{7pt} Three-Person Team \hspace{7pt}\textbullet\hspace{7pt} Spring 2026 \hspace{7pt}\textbullet\hspace{7pt} San Luis Obispo, California}

\node[anchor=north west,align=left,text width=10.00in,text=Ink,font=\fontsize{8.2}{9.5}\selectfont]
at ($(current page.north west)+(0.48in,-1.62in)$)
{This project asked a different perception question than object detection: instead of locating a sign with one bounding box, a \textbf{U-Net convolutional neural network} classified \textbf{every pixel} in the camera image as \texttt{line\_left}, \texttt{line\_right}, or \texttt{background}. That pixel-level output can preserve the shape and continuity of road markings, making it useful for downstream lane understanding.};

\node[anchor=north west,inner sep=1.7pt,draw=DeepGreen,line width=0.7pt]
at ($(current page.north west)+(0.48in,-2.28in)$)
{\includegraphics[width=6.05in,height=4.65in,keepaspectratio]{\detokenize{behind_shot.jpg}}};

\fill[DeepGreen] ($(current page.north west)+(0.65in,-2.50in)$) rectangle ($(current page.north west)+(2.73in,-2.23in)$);
\node[anchor=west,text=white,font=\sffamily\bfseries\fontsize{6.6}{7.2}\selectfont]
at ($(current page.north west)+(0.78in,-2.365in)$) {PHYSICAL CAMERA PLATFORM};

\node[anchor=north west,text=DeepGreen,font=\sffamily\bfseries\fontsize{8.4}{9.1}\selectfont]
at ($(current page.north west)+(6.48in,-2.28in)$) {ONE IMAGE BECOMES A PIXEL MAP};
\node[anchor=north west,align=left,text width=3.90in,text=Ink,font=\fontsize{7.5}{8.7}\selectfont]
at ($(current page.north west)+(6.48in,-2.66in)$)
{The training label is not a single category for the whole image. It is an image-sized mask in which each pixel carries one class ID. During inference, U-Net predicts the same kind of mask from a new camera frame.};

% Three semantic classes arranged in one horizontal row
\foreach \xx/\num/\ttl/\desc in {
6.48/01/LEFT LANE/\texttt{line\_left},
7.80/02/RIGHT LANE/\texttt{line\_right},
9.12/03/OTHER PIXELS/\texttt{background}}
{
\fill[SoftGreen,rounded corners=3pt]
($(current page.north west)+(\xx in,-3.38in)$)
rectangle
($(current page.north west)+(\xx in+1.22in,-4.08in)$);

\node[anchor=north west,text=PolyGold,font=\sffamily\bfseries\fontsize{6.4}{7.0}\selectfont]
at ($(current page.north west)+(\xx in+0.12in,-3.56in)$) {\num};

\node[anchor=north west,text=DeepGreen,font=\sffamily\bfseries\fontsize{6.5}{7.1}\selectfont]
at ($(current page.north west)+(\xx in+0.36in,-3.51in)$) {\ttl};

\node[anchor=north west,text=Muted,font=\sffamily\fontsize{5.8}{6.4}\selectfont]
at ($(current page.north west)+(\xx in+0.36in,-3.78in)$) {\desc};
}

\node[anchor=north west,inner sep=1.2pt,draw=RuleGray,line width=0.45pt]
at ($(current page.north west)+(6.48in,-4.14in)$)
{\includegraphics[width=3.94in]{\detokenize{intanceType.png}}};

\UNMetric{0.48in}{-7.28in}{2.77in}{LEFT-LANE PIXELS}{1.25\%}{of all labeled pixels}
\UNMetric{3.46in}{-7.28in}{2.77in}{RIGHT-LANE PIXELS}{0.78\%}{of all labeled pixels}
\UNMetric{6.44in}{-7.28in}{3.98in}{BACKGROUND PIXELS}{97.97\%}{expected majority in full-frame masks}
\UNFooter{\hyperlink{parallelparking}{\faArrowLeft\ PROJECT 13}}{1 / 6}{DATA QUALITY \faArrowRight}
\end{tikzpicture}

%==============================================================
%==============================================================
% PAGE 2 — BAD LABEL -> BAD STATISTICS -> BASELINE TRAINING
%==============================================================
\clearpage
\thispagestyle{empty}
\hypertarget{unet-data}{}
\null

\begin{tikzpicture}[remember picture,overlay]

\UNPageBase
\UNHeader{DATA QUALITY}{ONE BAD MASK CAN DISTORT AN ENTIRE DATASET}
{Annotation review
\hspace{7pt}\textbullet\hspace{7pt}
Pixel-count distributions
\hspace{7pt}\textbullet\hspace{7pt}
Baseline IoU + loss}

% -------------------------------------------------------------
% Narrative
% -------------------------------------------------------------
\node[
    anchor=north west,
    align=left,
    text width=10.00in,
    text=Ink,
    font=\fontsize{8.0}{9.3}\selectfont
]
at ($(current page.north west)+(0.48in,-1.62in)$)
{
Before changing the neural network, the dataset itself was audited.
A histogram of labeled pixels revealed \textbf{three extreme right-lane outliers}.
Inspection traced them to one incorrectly annotated frame and its augmented copies.
The page therefore follows a simple sequence:
\textbf{distribution anomaly $\rightarrow$ bad mask $\rightarrow$ baseline training behavior}.
};

% -------------------------------------------------------------
% TOP ROW — HISTOGRAM + BAD MASK
% -------------------------------------------------------------

% Left image: histogram
\node[
    anchor=north west,
    inner sep=1.2pt,
    draw=DeepGreen,
    line width=0.55pt
]
at ($(current page.north west)+(0.48in,-2.14in)$)
{
\includegraphics[
    width=4.78in,
    height=3.02in,
    keepaspectratio
]{\detokenize{histogram1.png}}
};

% Left banner over image
\fill[DeepGreen]
($(current page.north west)+(0.62in,-2.32in)$)
rectangle
($(current page.north west)+(2.30in,-2.04in)$);

\node[
    anchor=west,
    text=white,
    font=\sffamily\bfseries\fontsize{6.6}{7.1}\selectfont
]
at ($(current page.north west)+(0.76in,-2.18in)$)
{PIXEL OUTLIERS};

% Right image: bad annotation mask
\node[
    anchor=north west,
    inner sep=1.2pt,
    draw=DeepGreen,
    line width=0.55pt
]
at ($(current page.north west)+(5.56in,-2.14in)$)
{
\includegraphics[
    width=4.86in,
    height=3.02in,
    keepaspectratio
]{\detokenize{annotationMask1.png}}
};

% Right banner over image
\fill[DeepGreen]
($(current page.north west)+(5.70in,-2.32in)$)
rectangle
($(current page.north west)+(7.45in,-2.04in)$);

\node[
    anchor=west,
    text=white,
    font=\sffamily\bfseries\fontsize{6.6}{7.1}\selectfont
]
at ($(current page.north west)+(5.84in,-2.18in)$)
{BAD ANNOTATION};

% Short explanation directly beneath both images
\node[
    anchor=north west,
    align=left,
    text width=4.72in,
    text=Ink,
    font=\fontsize{7.0}{8.2}\selectfont
]
at ($(current page.north west)+(0.50in,-5.3in)$)
{
\textbf{\textcolor{DeepGreen}{What the histogram revealed.}}
Most right-lane masks occupied a compact pixel-count range, but three samples
contained dramatically more labeled pixels than the rest of the dataset.
};

\node[
    anchor=north west,
    align=left,
    text width=4.78in,
    text=Ink,
    font=\fontsize{7.0}{8.2}\selectfont
]
at ($(current page.north west)+(5.58in,-4.78in)$)
{
\textbf{\textcolor{DeepGreen}{Why those outliers existed.}}
One frame incorrectly labeled a large image region as \texttt{line\_right}.
Augmented copies inherited the same error, turning one annotation mistake into three outliers.
};

% -------------------------------------------------------------
% LOWER ROW — BASELINE TRAINING
% -------------------------------------------------------------
\draw[RuleGray,line width=0.5pt]
($(current page.north west)+(0.48in,-5.76in)$)
--
($(current page.north west)+(10.42in,-5.76in)$);

% Loss + IoU figure, moved up and enlarged
\node[
    anchor=north west,
    inner sep=1.2pt,
    draw=DeepGreen,
    line width=0.55pt
]
at ($(current page.north west)+(0.48in,-5.96in)$)
{
\includegraphics[
    width=6.82in,
    height=2.02in,
    keepaspectratio
]{\detokenize{lossANDiou1.png}}
};

% Banner over the plot
\fill[DeepGreen]
($(current page.north west)+(0.62in,-6.14in)$)
rectangle
($(current page.north west)+(2.52in,-5.86in)$);

\node[
    anchor=west,
    text=white,
    font=\sffamily\bfseries\fontsize{6.5}{7.0}\selectfont
]
at ($(current page.north west)+(0.76in,-6.00in)$)
{BASELINE TRAINING};


% Compact conclusion moved into the open right-side space
\node[
    anchor=north west,
    align=left,
    text width=2.86in,
    text=Ink,
    font=\fontsize{6.7}{7.8}\selectfont
]
at ($(current page.north west)+(7.55in,-5.90in)$)
{
\textbf{\textcolor{PolyGold}{WHY THIS MATTERED}}\\[3pt]
The model could still report reasonable aggregate metrics while learning from a flawed label distribution.
Cleaning the annotations therefore had to come \textbf{before} comparing architectures or hyperparameters.
};

% Metrics to the right
\UNMetric{7.55in}{-6.72in}{1.30in}{VAL IoU}{0.8628}{validation overlap}
\UNMetric{9.02in}{-6.72in}{1.40in}{TEST IoU}{0.8479}{unseen data}

% Compact training note below metrics, kept safely above footer
\node[
    anchor=north west,
    align=left,
    text width=2.87in,
    text=Ink,
    font=\fontsize{6.8}{7.8}\selectfont
]
at ($(current page.north west)+(7.55in,-7.51in)$)
{
\textbf{Training setup}\\
Batch 4, learning rate 0.001, maximum 100 epochs.
Early stopping ended the run at 34 epochs after validation loss stopped improving.
};


\UNFooter
{\faArrowLeft\ INTRODUCTION}
{2 / 6}
{CLEANED DATA \faArrowRight}

\end{tikzpicture}

% PAGE 3 — CLEANED DATA + BEST IOU RUN
%==============================================================
\clearpage
\thispagestyle{empty}
\hypertarget{unet-clean}{}
\null

\begin{tikzpicture}[remember picture,overlay]
\UNPageBase

\UNHeader
{CLEANED DATA}
{FIX THE LABELS FIRST, THEN JUDGE THE MODEL}
{Outlier removal
\hspace{7pt}\textbullet\hspace{7pt}
Cleaned class distribution
\hspace{7pt}\textbullet\hspace{7pt}
Controlled retraining}

% -------------------------------------------------------------
% PAGE NARRATIVE
% -------------------------------------------------------------
\node[
    anchor=north west,
    align=left,
    text width=10.00in,
    text=Ink,
    font=\fontsize{8.0}{9.3}\selectfont
]
at ($(current page.north west)+(0.48in,-1.62in)$)
{
The mislabeled frame and its augmented copies were removed before the next
training run. The purpose was not simply to improve the reported score, but to
restore a physically plausible right-lane distribution. The cleaned dataset was
then retrained with the \textbf{same batch size and learning rate}, allowing the
effect of correcting the labels to be evaluated directly.
};

% =============================================================
% STEP 1 — CLEANED PIXEL DISTRIBUTION
% =============================================================

\node[
    anchor=north west,
    text=DeepGreen,
    font=\sffamily\bfseries\fontsize{8.0}{8.7}\selectfont
]
at ($(current page.north west)+(0.48in,-2.10in)$)
{AFTER REMOVING THE OUTLIERS};

\node[
    anchor=north west,
    inner sep=1.2pt,
    draw=DeepGreen,
    line width=0.55pt
]
at ($(current page.north west)+(0.48in,-2.42in)$)
{
\includegraphics[
    width=5.25in,
    height=2.63in,
    keepaspectratio
]{\detokenize{histogram2.png}}
};

% Extended substantially leftward toward histogram; right edge remains fixed at 10.42in.
\fill[
    SoftGreen,
    rounded corners=4pt
]
($(current page.north west)+(4.48in,-2.42in)$)
rectangle
($(current page.north west)+(10.42in,-4.84in)$);

\node[
    anchor=north west,
    text=DeepGreen,
    font=\sffamily\bfseries\fontsize{8.0}{8.7}\selectfont
]
at ($(current page.north west)+(4.76in,-2.70in)$)
{WHAT CHANGED?};

\node[
    anchor=north west,
    align=left,
    text width=5.32in,
    text=Ink,
    font=\fontsize{7.3}{8.5}\selectfont
]
at ($(current page.north west)+(4.76in,-3.10in)$)
{
The three extreme right-lane samples are gone. Instead of extending to nearly
15,000 labeled pixels, the cleaned \texttt{line\_right} distribution now occupies
a range comparable to the left-lane samples.

\vspace{7pt}

\textbf{This changes what the experiment means.}
The network is no longer being evaluated against a label distribution distorted
by one annotation error and its augmented copies.

\vspace{7pt}

The next run therefore asks a controlled question:
\textbf{what happens when the same model is trained on cleaner data?}
};

% =============================================================
% DIVIDER
% =============================================================

\draw[
    RuleGray,
    line width=0.5pt
]
($(current page.north west)+(0.48in,-5.19in)$)
--
($(current page.north west)+(10.42in,-5.19in)$);

% =============================================================
% STEP 2 — RETRAIN ON CLEANED DATA
% =============================================================

\node[
    anchor=north west,
    text=DeepGreen,
    font=\sffamily\bfseries\fontsize{8.0}{8.7}\selectfont
]
at ($(current page.north west)+(0.48in,-5.48in)$)
{RETRAINING THE SAME MODEL ON CLEANER DATA};

\node[
    anchor=north west,
    inner sep=1.2pt,
    draw=RuleGray,
    line width=0.45pt
]
at ($(current page.north west)+(0.48in,-5.80in)$)
{
\includegraphics[
    width=6.62in,
    height=2.10in,
    keepaspectratio
]{\detokenize{lossANDiou2.png}}
};

% =============================================================
% CLEAN-DATA RESULTS
% =============================================================

\UNMetric
{7.40in}
{-5.80in}
{1.45in}
{VAL IoU}
{0.8714}
{highest of 4 runs}

\UNMetric
{8.97in}
{-5.80in}
{1.45in}
{TEST IoU}
{0.8508}
{highest of 4 runs}

\node[
    anchor=north west,
    align=left,
    text width=3.02in,
    text=Ink,
    font=\fontsize{6.9}{7.9}\selectfont
]
at ($(current page.north west)+(7.40in,-6.76in)$)
{
\textbf{Same basic training setup, cleaner labels.}

Batch 4, LR 0.001, with early stopping at 41 epochs.
This run produced the strongest validation and test IoU of the four experiments.

The remaining question was whether a good aggregate IoU also produced a
\textbf{visually trustworthy lane mask}.
};

\UNFooter
{\faArrowLeft\ DATA QUALITY}
{3 / 6}
{EXPERIMENTS \faArrowRight}

\end{tikzpicture}

%==============================================================
% PAGE 4 — WHY MORE TRAINING WAS TESTED
%==============================================================
\clearpage
\thispagestyle{empty}
\hypertarget{unet-experiments}{}
\null
\begin{tikzpicture}[remember picture,overlay]
\UNPageBase
\UNHeader{MODEL EXPERIMENTS}{WAS THE FALSE-POSITIVE ARTIFACT JUST OVERFITTING?}
{Batch-size experiment \hspace{7pt}\textbullet\hspace{7pt} Short vs. longer training \hspace{7pt}\textbullet\hspace{7pt} IoU + loss interpretation}

\node[anchor=north west,align=left,text width=10.00in,text=Ink,font=\fontsize{8.0}{9.3}\selectfont]
at ($(current page.north west)+(0.48in,-1.62in)$)
{A test prediction still contained small false-positive right-lane pixels. To investigate whether this came from excessive training, batch size was increased from 4 to 8 and two additional training lengths were examined. The key question was whether reducing training would remove the artifact without sacrificing useful lane segmentation.};

\node[anchor=north west,text=DeepGreen,font=\sffamily\bfseries\fontsize{8.1}{8.8}\selectfont]
at ($(current page.north west)+(0.48in,-2.38in)$) {RUN 3 — BATCH 8, 20 EPOCHS};
\node[anchor=north west,inner sep=1.2pt,draw=DeepGreen,line width=0.55pt]
at ($(current page.north west)+(0.48in,-2.72in)$)
{\includegraphics[width=4.82in,height=2.58in,keepaspectratio]{\detokenize{lossANDiou3.png}}};
\UNMetric{0.48in}{-4.64in}{2.22in}{TEST IoU}{0.8223}{model still improving}
\UNMetric{2.86in}{-4.64in}{2.44in}{TEST LOSS}{0.0183}{not yet converged}

\node[anchor=north west,text=DeepGreen,font=\sffamily\bfseries\fontsize{8.1}{8.8}\selectfont]
at ($(current page.north west)+(5.60in,-2.38in)$) {RUN 4 — BATCH 8, EARLY STOP AT 47 EPOCHS};
\node[anchor=north west,inner sep=1.2pt,draw=DeepGreen,line width=0.55pt]
at ($(current page.north west)+(5.60in,-2.72in)$)
{\includegraphics[width=4.82in,height=2.58in,keepaspectratio]{\detokenize{lossANDiou4.png}}};
\UNMetric{5.60in}{-4.64in}{2.22in}{TEST IoU}{0.8421}{recovered overlap}
\UNMetric{7.98in}{-4.64in}{2.44in}{TEST LOSS}{0.0159}{lowest clean-run test loss}

\fill[SoftGold,rounded corners=4pt] ($(current page.north west)+(0.48in,-5.48in)$) rectangle ($(current page.north west)+(10.42in,-7.31in)$);
\fill[PolyGold] ($(current page.north west)+(0.48in,-5.48in)$) rectangle ($(current page.north west)+(0.58in,-7.31in)$);
\node[anchor=north west,text=DeepGreen,font=\sffamily\bfseries\fontsize{8.0}{8.7}\selectfont]
at ($(current page.north west)+(0.78in,-6.02in)$) {WHAT THE EXPERIMENT SHOWED};
\node[anchor=north west,align=left,text width=9.18in,text=Ink,font=\fontsize{7.3}{8.5}\selectfont]
at ($(current page.north west)+(0.78in,-6.40in)$)
{The 20-epoch model had not finished improving, while the 47-epoch model recovered much of the IoU and reduced test loss. Because the same visible artifact persisted even when the training length changed, the evidence did \textbf{not} support the simple explanation that overfitting alone caused the false-positive lane fragments. Visual inspection still mattered alongside the metrics.};
% RESOURCE BAR — formatted consistently with the rest of the portfolio
% Replace the two placeholder URLs below when the final resources are ready.
\fill[SoftGreen!34,rounded corners=4pt]
($(current page.south west)+(0.46in,0.58in)$)
rectangle
($(current page.south east)+(-0.46in,1.08in)$);

\draw[RuleGray,line width=0.45pt,rounded corners=4pt]
($(current page.south west)+(0.46in,0.58in)$)
rectangle
($(current page.south east)+(-0.46in,1.08in)$);

\node[
anchor=center,
text=DeepGreen,
font=\sffamily\bfseries\fontsize{6.9}{7.7}\selectfont,
align=center
]
at ($(current page.south)+(0,0.83in)$)
{
\href{https://example.com/project-report.pdf}
{\faFilePdf\ PROJECT REPORT}
\hspace{18pt}
\href{https://colab.research.google.com/}
{\faCode\ COLAB / GITHUB}
};

\UNFooter
{\faArrowLeft\ CLEANED DATA}
{4 / 6}
{PREDICTIONS \faArrowRight}
\end{tikzpicture}

%==============================================================
% PAGE 5 — SAME SCENE, THREE MODEL PREDICTIONS
%==============================================================
\clearpage
\thispagestyle{empty}
\hypertarget{unet-predictions}{}
\null

\begin{tikzpicture}[remember picture,overlay]
\UNPageBase

\UNHeader
{PREDICTIONS}
{WHAT THE NETWORK ACTUALLY PAINTED}
{Same test scene
\hspace{7pt}\textbullet\hspace{7pt}
Ground truth vs. prediction
\hspace{7pt}\textbullet\hspace{7pt}
Qualitative comparison}

% -------------------------------------------------------------
% INTRODUCTION
% -------------------------------------------------------------
\node[
    anchor=north west,
    align=left,
    text width=10.00in,
    text=Ink,
    font=\fontsize{8.0}{9.3}\selectfont
]
at ($(current page.north west)+(0.48in,-1.62in)$)
{
IoU and loss summarize millions of pixel decisions into single numbers.
These prediction images show what those numbers can hide: the dominant lane
structure is usually recovered, while small right-lane fragments can appear
near reflective or high-contrast regions. Viewing the same scene across the
three runs makes the progression easier to judge.
};

% =============================================================
% LEFT COLUMN — THREE PREDICTION IMAGES
% =============================================================

% ----------------------------- RUN 1 ---------------------------
\node[
    anchor=north west,
    inner sep=1.0pt,
    draw=DeepGreen,
    line width=0.55pt
]
at ($(current page.north west)+(0.52in,-2.40in)$)
{
\includegraphics[
    width=6.15in,
    height=1.68in,
    keepaspectratio
]{\detokenize{annotationMask2.png}}
};

% compact overlay banner
\fill[DeepGreen,rounded corners=1.5pt]
($(current page.north west)+(0.66in,-2.56in)$)
rectangle
($(current page.north west)+(2.46in,-2.36in)$);

\node[
    anchor=west,
    text=white,
    font=\sffamily\bfseries\fontsize{5.9}{6.4}\selectfont
]
at ($(current page.north west)+(0.80in,-2.46in)$)
{RUN 1 \hspace{5pt} UNMODIFIED DATA};

% ----------------------------- RUN 2 ---------------------------
\node[
    anchor=north west,
    inner sep=1.0pt,
    draw=DeepGreen,
    line width=0.55pt
]
at ($(current page.north west)+(0.52in,-4.18in)$)
{
\includegraphics[
    width=6.15in,
    height=1.68in,
    keepaspectratio
]{\detokenize{annotationMask3.png}}
};

\fill[DeepGreen,rounded corners=1.5pt]
($(current page.north west)+(0.66in,-4.34in)$)
rectangle
($(current page.north west)+(2.93in,-4.14in)$);

\node[
    anchor=west,
    text=white,
    font=\sffamily\bfseries\fontsize{5.9}{6.4}\selectfont
]
at ($(current page.north west)+(0.80in,-4.24in)$)
{RUN 2 \hspace{5pt} CLEAN DATA / BEST IoU};

% ----------------------------- RUN 3 ---------------------------
\node[
    anchor=north west,
    inner sep=1.0pt,
    draw=DeepGreen,
    line width=0.55pt
]
at ($(current page.north west)+(0.52in,-5.96in)$)
{
\includegraphics[
    width=6.15in,
    height=1.68in,
    keepaspectratio
]{\detokenize{annotationMask4.png}}
};

\fill[DeepGreen,rounded corners=1.5pt]
($(current page.north west)+(0.66in,-6.12in)$)
rectangle
($(current page.north west)+(2.73in,-5.92in)$);

\node[
    anchor=west,
    text=white,
    font=\sffamily\bfseries\fontsize{5.9}{6.4}\selectfont
]
at ($(current page.north west)+(0.80in,-6.02in)$)
{RUN 3 \hspace{5pt} SHORT BATCH-8 TEST};

% =============================================================
% RIGHT COLUMN — INTERPRETATION
% =============================================================

\node[
    anchor=north west,
    text=DeepGreen,
    font=\sffamily\bfseries\fontsize{8.0}{8.7}\selectfont
]
at ($(current page.north west)+(5.79in,-2.42in)$)
{HOW TO READ THESE IMAGES};

\node[
    anchor=north west,
    align=left,
    text width=4.41in,
    text=Ink,
    font=\fontsize{7.2}{8.4}\selectfont
]
at ($(current page.north west)+(5.79in,-2.78in)$)
{
The cleaned 41-epoch model gives the strongest lane overlap numerically,
but the shorter batch-8 experiment does not eliminate the false-positive
artifact.

\vspace{6pt}

The report attributes part of the erroneous right-lane prediction to glare
in the scene. That makes the qualitative comparison essential:
\textbf{good segmentation is not only a high score, but a physically sensible mask.}
};

% -------------------------------------------------------------
% WHAT IoU MEANS HERE — same-width right-column panel
% -------------------------------------------------------------
\fill[SoftGreen,rounded corners=4pt]
($(current page.north west)+(5.71in,-4.55in)$)
rectangle
($(current page.north west)+(10.42in,-7.64in)$);

\node[
    anchor=north west,
    text=DeepGreen,
    font=\sffamily\bfseries\fontsize{7.8}{8.5}\selectfont
]
at ($(current page.north west)+(5.99in,-4.83in)$)
{WHAT IoU MEANS HERE};

\node[
    anchor=north west,
    align=left,
    text width=4.08in,
    text=Ink,
    font=\fontsize{6.75}{7.75}\selectfont
]
at ($(current page.north west)+(5.99in,-5.20in)$)
{
\textbf{Intersection over Union (IoU)} measures the overlap between a predicted
segmentation mask and the labeled ground truth. A higher IoU means the model is
assigning more pixels to the correct class while avoiding unnecessary
false-positive regions.

\vspace{6pt}

For semantic segmentation, IoU is more informative than raw accuracy because
the \textbf{background class dominates most of the image}. A model could achieve
high pixel accuracy by predicting background nearly everywhere while still
missing the lane markings that matter for autonomy.

\vspace{6pt}

\textbf{Training loss} describes how strongly the network is penalized during
optimization, while \textbf{validation loss} checks whether that learned behavior
generalizes to data it did not directly train on. When training improves but
validation performance stops improving, the model may be beginning to
\textbf{overfit}.

\vspace{6pt}

These images also show why quantitative metrics need a qualitative check.
Two runs can produce similar IoU values while making very different mistakes.
A small false-positive fragment can be numerically minor but still physically
misleading if it appears where no lane exists.

\vspace{6pt}

\textbf{Key machine-learning ideas shown here:}
class imbalance, semantic segmentation, train/validation behavior,
generalization, early stopping, overfitting, and qualitative model validation.
};

% =============================================================
% FOOTER
% =============================================================
\UNFooter
{\faArrowLeft\ EXPERIMENTS}
{5 / 6}
{RESULTS \faArrowRight}

\end{tikzpicture}

%==============================================================
% PAGE 6 — SIMPLE RESULTS + PHYSICAL CONTEXT + RESOURCES
%==============================================================
\clearpage
\thispagestyle{empty}
\hypertarget{unet-results}{}
\null
\begin{tikzpicture}[remember picture,overlay]
\UNPageBase
\UNHeader{PROJECT RESULT}{FROM LABELED PIXELS TO A ROAD-SCENE PERCEPTION MODEL}
{Dataset diagnosis \hspace{7pt}\textbullet\hspace{7pt} U-Net training \hspace{7pt}\textbullet\hspace{7pt} Physical lane-perception context}

\node[anchor=north west,align=left,text width=10.00in,text=Ink,font=\fontsize{8.0}{9.3}\selectfont]
at ($(current page.north west)+(0.48in,-1.62in)$)
{The final comparison did not produce one universally ``best'' model. The clean-data batch-4 run achieved the strongest IoU, while the longer batch-8 run achieved the lowest test loss among the cleaned-data experiments and reduced the visible artifact. The useful conclusion was therefore a \textbf{model-selection tradeoff}, not simply a ranking.};

% simplified table-style results
\fill[SoftGreen,rounded corners=5pt] ($(current page.north west)+(0.48in,-2.34in)$) rectangle ($(current page.north west)+(6.56in,-5.84in)$);
\fill[DeepGreen,rounded corners=5pt] ($(current page.north west)+(0.48in,-2.34in)$) rectangle ($(current page.north west)+(6.56in,-2.78in)$);
\node[anchor=west,text=white,font=\sffamily\bfseries\fontsize{7.0}{7.6}\selectfont]
at ($(current page.north west)+(0.76in,-2.56in)$) {EXPERIMENT};
\node[anchor=center,text=white,font=\sffamily\bfseries\fontsize{7.0}{7.6}\selectfont]
at ($(current page.north west)+(4.33in,-2.56in)$) {TEST IoU};
\node[anchor=east,text=white,font=\sffamily\bfseries\fontsize{7.0}{7.6}\selectfont]
at ($(current page.north west)+(6.25in,-2.56in)$) {TEST LOSS};

\foreach \yy/\model/\iou/\loss/\tag in {
3.10/Unmodified · B4 · 34 ep/0.8479/0.0165/BASELINE,
3.72/Clean · B4 · 41 ep/0.8508/0.0186/BEST IoU,
4.34/Clean · B8 · 20 ep/0.8223/0.0183/UNDERTRAINED,
4.96/Clean · B8 · 47 ep/0.8421/0.0159/LOWEST CLEAN LOSS}
{
\node[anchor=west,text=Ink,font=\fontsize{6.9}{7.6}\selectfont] at ($(current page.north west)+(0.76in,-\yy in)$) {\model};
\node[anchor=center,text=Ink,font=\bfseries\fontsize{7.0}{7.6}\selectfont] at ($(current page.north west)+(4.33in,-\yy in)$) {\iou};
\node[anchor=east,text=Ink,font=\bfseries\fontsize{7.0}{7.6}\selectfont] at ($(current page.north west)+(6.25in,-\yy in)$) {\loss};
\node[anchor=west,text=DeepGreen,font=\sffamily\bfseries\fontsize{5.8}{6.3}\selectfont] at ($(current page.north west)+(0.76in,-\yy in-0.23in)$) {\tag};
}

\node[anchor=north west,text=DeepGreen,font=\sffamily\bfseries\fontsize{8.0}{8.7}\selectfont]
at ($(current page.north west)+(6.86in,-2.34in)$) {PHYSICAL CONTEXT};
\node[anchor=north west,inner sep=1.4pt,draw=DeepGreen,line width=0.55pt]
at ($(current page.north west)+(6.86in,-2.68in)$)
{\includegraphics[width=3.56in,height=2.72in,keepaspectratio]{\detokenize{physical_framework.jpg}}};
\node[anchor=north west,align=left,text width=3.56in,text=Ink,font=\fontsize{6.9}{8.0}\selectfont]
at ($(current page.north west)+(6.86in,-5.50in)$)
{The segmentation study was developed for the same physical autonomous-vehicle environment: camera imagery becomes a lane mask that can support higher-level road interpretation and control.};

\draw[RuleGray,line width=0.5pt] ($(current page.north west)+(0.48in,-6.28in)$) -- ($(current page.north west)+(10.42in,-6.28in)$);
\node[anchor=north west,text=DeepGreen,font=\sffamily\bfseries\fontsize{8.1}{8.8}\selectfont]
at ($(current page.north west)+(0.48in,-6.55in)$) {WHAT THIS PROJECT DEMONSTRATED};

\foreach \xx/\ttl/\txt in {
0.48/DATA QUALITY/Histogram analysis exposed a labeling failure before model metrics hid it.,
3.78/MODEL EVALUATION/IoU + loss were compared with visual prediction quality rather than used alone.,
7.08/ENGINEERING WORKFLOW/Clean labels → train → diagnose predictions → adjust experiment → compare again.}
{
\fill[SoftGold,rounded corners=3pt] ($(current page.north west)+(\xx in,-6.92in)$) rectangle ++(3.02in,-1.03in);
\node[anchor=north west,text=DeepGreen,font=\sffamily\bfseries\fontsize{7.0}{7.6}\selectfont] at ($(current page.north west)+(\xx in+0.18in,-7.12in)$) {\ttl};
\node[anchor=north west,align=left,text width=2.64in,text=Ink,font=\fontsize{6.5}{7.4}\selectfont] at ($(current page.north west)+(\xx in+0.18in,-7.43in)$) {\txt};
}

\node[anchor=north west,text=DeepGreen,font=\sffamily\bfseries\fontsize{7.0}{7.6}\selectfont]
at ($(current page.north west)+(0.50in,-8.18in)$)
{\faLaptopCode\ \ U-NET COLAB NOTEBOOK \hspace{0.55in} \faFilePdf\ \ PROJECT REPORT};

\UNFooter{\faArrowLeft\ PREDICTIONS}{6 / 6}{PROJECT 15 \faArrowRight}
\end{tikzpicture}


\clearpage
